% !TeX root = ../../Book3.tex
\section{Hurwitz 定理与全纯极限的应用}
\label{b3:sec:0404}

局部一致收敛已经保证极限全纯，却还没有回答极限能否取得某个值、是否仍为单射。第三章的零点计数提供了进一步的控制：在一个无零边界上，极限函数的模有正下界；函数列与它的误差最终小于这个下界，Rouché 定理便把内部零点的总重数传回函数列。本节先沿这条路线研究零点和映射，再把 Montel 紧性与恒等定理结合，说明何时一组点值极限足以决定整列的局部一致极限。

\subsection{零点稳定性}
\label{b3:subsec:040401}

零点的位置可以变化，几个单零点也可以在极限中合并。稳定的量是一个固定围道内部、按重数计算的零点总数，而保证稳定的条件是极限函数在围道上不为零。\term[hurwitz-dingli]{Hurwitz 定理}[Hurwitz's theorem]把这一观察用于全纯函数列。

\Needspace{6\baselineskip}
\begin{theorem}[Hurwitz]\label{b3:thm:04-hurwitz}
设 \(U\subseteq\mathbb C\) 为区域，\(f_n:U\rightarrow\mathbb C\) 全纯，并在 \(U\) 上局部一致收敛于有限值函数 \(f\)，且 \(f\) 不恒为零。若 \(a\in U\)、\(r>0\)、\(\overline{D(a,r)}\subseteq U\)，且 \(f\) 在 \(|z-a|=r\) 上无零点，则充分大的 \(n\) 使 \(f_n\) 在该圆周无零点，并与 \(f\) 在 \(D(a,r)\) 中具有相同的零点总重数。

\ProseParagraphBreak

此外，对每个满足 \(f(a)=0\) 的点 \(a\in U\)，存在整数 \(N\ge1\) 及点 \(a_n\in U\)、\(n\ge N\)，使 \(f_n(a_n)=0\) 且 \(a_n\rightarrow a\)。
\end{theorem}
\begin{proof}
由\cref{b3:thm:04-weierstrass}，\(f\) 全纯。它不恒零，所以零点孤立，在所给闭圆盘中只有有限个零点。圆周的紧性与无零条件给出
\[
 \delta=\min_{|z-a|=r}|f(z)|>0.
\]
局部一致收敛使充分大的 \(n\) 满足
\[
 \max_{|z-a|=r}|f_n(z)-f(z)|<\delta\le |f(z)|
 \quad(|z-a|=r).
\]
于是 \(f_n\) 在圆周也无零点。对全纯主项 \(f\) 和余项 \(f_n-f\) 应用\cref{b3:thm:03-rouche}，便得圆盘内的零点总重数相同。

再设 \(a\) 是 \(f\) 的 \(m\) 阶零点。取严格递减到零的半径 \(r_j>0\)，使各闭圆盘均包含于 \(U\)，且其中 \(f\) 的唯一零点为 \(a\)。由刚证明的计数结论，对每个 \(j\) 有整数 \(M_j\)，使全部 \(n\ge M_j\) 的 \(f_n\) 在 \(D(a,r_j)\) 中均有 \(m\) 个零点，计重数。将这些最终指标增大，递归取
\[
 N_1\ge M_1,\quad
 N_{j+1}\ge\max\{M_{j+1},N_j+1\}.
\]
对每个 \(n\ge N_1\)，令
\[
 j(n)=\max\{\,j\ge1\mid N_j\le n\,\}.
\]
由于 \(N_j\) 严格增加，这个最大值存在且有限。在 \(D(a,r_{j(n)})\) 中选一个 \(f_n\) 的零点 \(a_n\)。当 \(n\rightarrow\infty\) 时，\(j(n)\rightarrow\infty\)，故 \(|a_n-a|<r_{j(n)}\rightarrow0\)。这一选择对每个足够大的指标 \(n\) 都有定义，给出了所需的整条零点列。
\end{proof}

最后这一步包含两个不同的尺度。先固定半径，才能找到适用于该圆盘的最终指标；随后让半径按层缩小，并相应推迟指标，才能把“每个小盘最终有根”写成 \(a_n\rightarrow a\)。定理不要求给同一重根分裂出的各个根作唯一编号。比如 \(z^2-1/n\) 的两个单零点 \(\pm n^{-1/2}\) 都趋于零，而极限 \(z^2\) 在零点处的重数为二。

\ProseParagraphBreak

作为一个不能直接用多项式求根公式处理的例子，考虑
\[
 f_n(z)=z^2-1+\frac{e^z}{n}.
\]
在每个固定紧集上，末项一致趋于零，所以极限为 \(z^2-1\)。对 \(a=1\) 或 \(a=-1\)，在 \(|z-a|=1/4\) 上都有
\[
 |z^2-1|=|z-a|\cdot|z+a|\ge\frac14\cdot\frac74=\frac7{16},
 \quad \frac{|e^z|}{n}\le\frac{e^{5/4}}n.
\]
因此只要 \(n>16e^{5/4}/7\)，每个圆盘 \(D(1,1/4)\)、\(D(-1,1/4)\) 中都恰有一个计重数零点，因而都是单零点。继续缩小这两个圆盘，Hurwitz 定理说明相应零点分别趋于 \(1\) 和 \(-1\)。这里的边界下界既说明为何能数根，也给出了固定圆盘适用的具体指标。

\ProseParagraphBreak

反过来，若 \(f_n(a_n)=0\)，各点 \(a_n\) 始终落在某个紧集 \(K\subseteq U\) 中，且 \(a_n\rightarrow a\in K\)，则
\[
 |f(a)|\le |f(a)-f(a_n)|+\sup_K|f-f_n|\longrightarrow0,
\]
所以聚点也是极限函数的零点。共同紧集不能省略：\(f_n(z)=1-z/n\) 在全平面局部一致趋于 \(1\)，但它的零点 \(n\) 逃向无穷远。Hurwitz 的结论只约束固定内部区域中的零点。

\ProseParagraphBreak

若极限恒零，计数便失去上述依据。函数列 \(1/n\) 无零，函数列 \(z/n\) 每个成员在零点处有一个单零点，两列却都局部一致趋于零。恒零函数没有离散零点集，也无法在任何圆周上提供严格正的最低模；不能把它看成具有某个有限总重数的函数继续套用计数结论。

\subsection{无零极限}
\label{b3:subsec:040402}

没有零点是一种在极限下几乎封闭的性质。唯一需要留下的退化情形，是整个函数同时缩小到零。

\Needspace{6\baselineskip}
\begin{corollary}[无零极限]\label{b3:cor:04-zero-free-limit}
设 \(U\subseteq\mathbb C\) 为区域，\(f_n:U\rightarrow\mathbb C\) 全纯，且对全部 \(n\ge1\)、\(z\in U\) 有 \(f_n(z)\ne0\)。若 \(f_n\) 在 \(U\) 上局部一致收敛于有限值函数 \(f\)，则 \(f\) 全纯，并且或者在 \(U\) 恒为零，或者在 \(U\) 处处无零。
\end{corollary}
\begin{proof}
全纯性由 Weierstrass 定理给出。若 \(f\) 不恒零却在 \(a\in U\) 为零，取一个闭包包含于 \(U\) 的小圆盘，使 \(a\) 是其中唯一的零点。其重数为某个正整数，边界无零。Hurwitz 定理迫使全部充分晚的 \(f_n\) 在该圆盘中也有零点，与假设矛盾。
\end{proof}

常值函数列 \(1/n\) 说明恒零例外确会发生。若 \(U\) 上的全纯函数列 \(f_n\) 共同省略某个固定值 \(b\in\mathbb C\)，且局部一致趋于有限值 \(f\)，则对 \(f_n-b\) 应用推论：\(f\) 要么恒为 \(b\)，要么仍处处不取值 \(b\)。平移到零点以后，省略值问题也变成了零点稳定性。

\ProseParagraphBreak

当极限仍无零时，倒数也有通常的局部一致极限。具体地，设 \(U\) 上的无零全纯函数 \(f_n\) 局部一致趋于处处非零的 \(f\)。固定非空紧集 \(K\subseteq U\)，有 \(c_K=\min_K|f|>0\)，且充分大的 \(n\) 使 \(|f_n|\ge c_K/2\) 于 \(K\)。于是
\[
 \sup_K\left|\frac1{f_n}-\frac1f\right|
 \le\frac2{c_K^2}\cdot\sup_K|f_n-f|\longrightarrow0.
\]
分母的共同正下界是取倒数与极限交换的依据；若 \(f_n=1/n\)，倒数则趋于无穷，不能得到有限值版本的结论。

\ProseParagraphBreak

区域的连通性同样有作用。令 \(U=D(-2,1)\cup D(2,1)\)，在左边的圆盘取 \(f_n=1/n\)，右边取 \(f_n=1\)。每个成员在 \(U\) 全纯且无零，极限却在一个分支恒零、在另一个分支无零。对一般开集，无零极限推论应逐连通分支使用。若考虑全纯函数列的球面局部一致极限，\cref{b3:cor:04-holomorphic-spherical-limit}已说明它只能是全纯函数或全域恒为无穷，不会在内点新生孤立极点。

\subsection{单叶极限}
\label{b3:subsec:040403}

区域上的单射全纯函数称为\term[danye-hanshu]{单叶函数}[univalent function]。它的不同点不能映到同一值。要研究这一性质在极限中是否保留，可以固定一个点 \(a\)，把“另一个点也映到 \(f(a)\)”改写成差商的零点问题；在中心补上导数，便得到可使用无零极限推论的全纯函数。

\Needspace{6\baselineskip}
\begin{theorem}[单叶极限]\label{b3:thm:04-univalent-limit}
设 \(U\subseteq\mathbb C\) 为区域，\(f_n:U\rightarrow\mathbb C\) 是单射全纯函数，并在 \(U\) 上局部一致收敛于有限值函数 \(f\)。则 \(f\) 是常值函数，或者是单射全纯函数。
\end{theorem}
\begin{proof}
Weierstrass 定理保证 \(f\) 全纯。假设 \(f\) 非常值，任取 \(a\in U\)，在 \(z\ne a\) 处定义
\[
 g_n(z)=\frac{f_n(z)-f_n(a)}{z-a},
 \quad g(z)=\frac{f(z)-f(a)}{z-a},
\]
并令 \(g_n(a)=f_n'(a)\)、\(g(a)=f'(a)\)。Taylor 展开说明补定义后的函数在整个 \(U\) 全纯。对 \(z\ne a\)，\(f_n\) 的单射性给出 \(g_n(z)\ne0\)；在中心处，\cref{b3:cor:01-local-injectivity}说明单射全纯函数的导数不为零，故也有 \(g_n(a)\ne0\)。

须先核对 \(g_n\rightarrow g\) 是局部一致收敛。取 \(\rho>0\)，使 \(\overline{D(a,\rho)}\subseteq U\)。在此闭圆盘上，沿直线段积分给出
\[
 g_n(z)=\int_0^1 f_n'\bigl(a+t(z-a)\bigr)\,\mathrm{d}t,
\]
\(g\) 也有对应表达。所有线段都留在同一个闭圆盘中，故
\[
 \sup_{\overline{D(a,\rho)}}|g_n-g|
 \le\sup_{\overline{D(a,\rho)}}|f_n'-f'|\longrightarrow0,
\]
其中用了\cref{b3:prop:04-derivative-convergence}。对任意紧集 \(K\subseteq U\)，其在 \(D(a,\rho/2)\) 外的部分满足 \(|z-a|\ge\rho/2\)，所以
\[
 |g_n(z)-g(z)|
 \le\frac2\rho\bigl(|f_n(z)-f(z)|+|f_n(a)-f(a)|\bigr).
\]
这部分也一致收敛，与中心附近的估计合起来即得全 \(K\) 上的一致收敛。

由于 \(f\) 非常值，\(g\) 不恒为零。对无零函数列 \(g_n\) 使用\cref{b3:cor:04-zero-free-limit}，得 \(g\) 处处无零。因而对每个 \(z\ne a\) 都有 \(f(z)\ne f(a)\)。点 \(a\) 任意，所以 \(f\) 单射。
\end{proof}

函数列 \(f_n(z)=z/n\) 在任何固定区域上都单叶，却局部一致趋于常值零。若在某个固定点 \(a\in U\) 另有 \(|f_n'(a)|\ge c>0\)，则导数收敛给出 \(|f'(a)|\ge c\)，从而排除常值退化。这里只固定 \(f_n(a)\) 的值还不够，\(z/n\) 在原点已经具有相同的函数值。

\ProseParagraphBreak

映射的收敛还提出了一个定义域问题。即使 \(f_n(z)\rightarrow f(z)\)，给定 \(w\in f(U)\) 后，方程 \(f_n(z)=w\) 是否有解仍需证明；只有先确认这一点，才谈得上 \(f_n^{-1}(w)\) 的极限。单叶全纯函数导数无零，由局部逆函数定理，其逆函数在像区域上全纯。下面证明极限像的每个紧集最终都落在这些逆函数的定义域中。

\Needspace{6\baselineskip}
\begin{proposition}[单叶逆映射的收敛]\label{b3:prop:04-inverse-convergence}
设 \(U\subseteq\mathbb C\) 为区域，\(f_n:U\rightarrow\mathbb C\) 为单叶全纯函数，并在 \(U\) 上局部一致收敛于非常值函数 \(f\)。则 \(f\) 单叶全纯，且对每个紧集 \(K\subseteq f(U)\)，存在整数 \(N\)，使全部 \(n\ge N\) 满足 \(K\subseteq f_n(U)\)，并有
\[
 \sup_{w\in K}|f_n^{-1}(w)-f^{-1}(w)|\longrightarrow0.
\]
\end{proposition}
\begin{proof}
单叶性由前一定理保证。先固定 \(w_0\in f(U)\)，记 \(a=f^{-1}(w_0)\)。取 \(r>0\)，使闭圆盘 \(C=\overline{D(a,r)}\) 包含于 \(U\)。由于 \(f\) 单射，\(f-w_0\) 在该圆盘中只有一个单零点 \(a\)，边界不为零。因此
\[
 m=\min_{|z-a|=r}|f(z)-w_0|>0.
\]
先让目标点在 \(D(w_0,m/3)\) 中移动。对其中的每个 \(w\)，常数扰动 \(|w-w_0|<m/3\) 小于圆周上的 \(|f-w_0|\)。Rouché 定理说明 \(f-w\) 在源圆盘中也恰有一个计重数零点；同时圆周上有共同下界
\[
 |f(z)-w|\ge |f(z)-w_0|-|w-w_0|>2m/3.
\]
局部一致收敛提供一个与 \(w\) 无关的最终指标，使圆周上 \(|f_n-f|<m/3\)。第二次应用 Rouché 定理，比较 \(f-w\) 与 \(f_n-w\)，得 \(f_n-w\) 在该圆盘中同样恰有一个零点。故整个 \(D(w_0,m/3)\) 同时包含于 \(f(U)\) 和这些 \(f_n(U)\)，且相应逆像都落在 \(D(a,r)\) 中。

现在在较小闭盘 \(K_0=\overline{D(w_0,m/4)}\) 上比较逆像。对 \(w\in K_0\)，记 \(z_n=f_n^{-1}(w)\)、\(z=f^{-1}(w)\)。它们都在紧集 \(C\) 内，且
\[
 |f(z_n)-f(z)|=|f(z_n)-f_n(z_n)|
 \le\sup_C|f-f_n|\longrightarrow0.
\]
连续单射 \(f:C\rightarrow f(C)\) 是从紧空间到 Hausdorff 空间的同胚，故其逆映射在紧集 \(f(C)\) 上一致连续。这里 \(f(z_n)\) 和 \(f(z)\) 均属于 \(f(C)\)，所以给定 \(\varepsilon>0\)，可选一个对全部 \(w\in K_0\) 有效的 \(\delta>0\)，使上式左端小于 \(\delta\) 时便有 \(|z_n-z|<\varepsilon\)。因而逆映射在 \(K_0\) 上一致收敛。

对给定紧集 \(K\subseteq f(U)\) 的每个点执行这一构造，所得较小像盘的内部覆盖 \(K\)。取有限子覆盖，并将各处的包含指标及一致误差指标取最大，便同时得到 \(K\subseteq f_n(U)\) 和整个 \(K\) 上的逆映射一致收敛。
\end{proof}

逆映射的定义域 \(f_n(U)\) 随指标变化，所以这里的收敛按紧集理解：先固定极限像中的一个紧集，再从某项起把所有逆映射限制到这个共同集合上。Rouché 比较中的边界正距离同时控制了目标点 \(w\) 的移动和近似函数 \(f_n\) 的误差，才使最终指标对整块像盘统一有效。

\ProseParagraphBreak

例如在单位圆盘 \(\mathbb D\) 上取
\[
 f_n(z)=\frac{z}{1+z/n}.
\]
分母在盘内无零，交叉相乘可核对 \(f_n\) 单射；在任意 \(|z|\le r<1\) 上，\(|f_n(z)-z|\le r^2/(n-r)\)，所以局部一致趋于恒等映射。解方程得到逆函数候选式
\[
 f_n^{-1}(w)=\frac{w}{1-w/n}.
\]
对 \(|w|\le r<1\)，当 \(n>r/(1-r)\) 时，这个式子的模至多为 \(r/(1-r/n)<1\)，确实给出位于原定义域中的逆像，而且
\[
 |f_n^{-1}(w)-w|\le\frac{r^2}{n(1-r/n)}\longrightarrow0.
\]
但对每个固定 \(n\)，选实数 \(n/(n+1)<w<1\)，同一逆函数候选值却大于 \(1\)，所以该 \(w\) 不属于 \(f_n(\mathbb D)\)。因此每个内部紧集最终被覆盖，并不意味着整个极限像 \(\mathbb D\) 最终都包含于 \(f_n(\mathbb D)\)。

\subsection{Vitali 型收敛判据}
\label{b3:subsec:040404}

Montel 定理提供收敛子列，还不能单独确定它们是否具有同一个极限。恒等定理恰好提供所需的唯一性：两个全纯极限只要在具有内部聚点的集合上相等，就在整个区域相等。把这两项结合，得到下面的\term[vitali-shoulian-dingli]{Vitali 收敛定理}[Vitali convergence theorem]的局部有界形式。

\Needspace{6\baselineskip}
\begin{proposition}[Vitali，局部有界形式]\label{b3:prop:04-vitali-bounded}
设 \(U\subseteq\mathbb C\) 为区域，\(f_n:U\rightarrow\mathbb C\) 全纯，并且对每个紧集 \(K\subseteq U\)，存在有限常数 \(M_K\)，使全部 \(n\ge1\) 满足 \(\sup_K|f_n|\le M_K\)。设 \(E\subseteq U\) 在 \(U\) 内有聚点，且对每个 \(z\in E\)，数列 \(f_n(z)\) 收敛于有限值。则存在全纯函数 \(f:U\rightarrow\mathbb C\)，使整列 \(f_n\) 在 \(U\) 上局部一致收敛于 \(f\)。
\end{proposition}
\begin{proof}
局部有界性使\cref{b3:thm:04-montel}适用，先从原序列抽取一条局部一致收敛子列，记其全纯极限为 \(f\)。这一步给出了候选极限的存在。

再任取原序列的一条子列。它仍属于同一个局部有界族，故 Montel 定理使它有局部一致收敛的子子列，记全纯极限为 \(g\)。对每个 \(z\in E\)，两条已收敛的子列都必须趋于原数列 \(f_n(z)\) 的点值极限，所以 \(f(z)=g(z)\)。集合 \(E\) 有位于 \(U\) 内部的聚点，区域 \(U\) 又连通，\cref{b3:thm:01-identity}因此给出 \(g=f\) 于整个 \(U\)。于是原序列的每条子列都能再抽出以 \(f\) 为局部一致极限的子子列。

最后证明整列收敛。若在某个紧集 \(K\subseteq U\) 上不一致趋于 \(f\)，则存在 \(\varepsilon>0\) 及严格增加的指标 \(n_j\)，使
\[
 \sup_K|f_{n_j}-f|\ge\varepsilon\quad(j\ge1).
\]
由前一步，这条子列又有局部一致趋于 \(f\) 的子子列，因而在 \(K\) 上的误差趋于零，与固定正下界矛盾。
\end{proof}

这一判据尤其适合从实变量极限获得复变量结论。令
\[
 F_n(z)=\left(1+\frac zn\right)^n\quad(z\in\mathbb C).
\]
每个 \(F_n\) 都是整函数，且
\[
 |F_n(z)|\le\left(1+\frac{|z|}{n}\right)^n\le e^{|z|}.
\]
任意固定紧集上的 \(|z|\) 有上界，故这给出了全族的局部共同界。在实区间 \((-1,1)\) 上，实变量的指数极限给出 \(F_n(x)\rightarrow e^x\)。Vitali 判据使整列在全平面局部一致趋于某个整函数 \(F\)；\(F\) 与 \(e^z\) 在这个实区间相等，恒等定理于是确定 \(F=e^z\)。

\ProseParagraphBreak

得到局部一致收敛后，才可以应用 Weierstrass 定理取导数。对每个固定整数 \(k\ge1\)，在 \(n\ge k\) 时有
\[
 F_n^{(k)}(z)=\frac{n(n-1)\cdots(n-k+1)}{n^k}
 \cdot\left(1+\frac zn\right)^{n-k}\longrightarrow e^z,
\]
收敛在每个紧集上一致成立。这条推导只在实区间使用了最初的点值极限，区域中的共同上界和全纯唯一性负责把它扩展到整个复平面。

\ProseParagraphBreak

局部有界条件不能删除。在 \(\mathbb D\) 上取 \(f_n(z)=(2z)^n\)，则对实区间 \(E=(-1/2,1/2)\) 中每个点都有 \(f_n(x)\rightarrow0\)，而 \(E\) 具有内部聚点；可是 \(f_n(3/4)=(3/2)^n\) 发散。该序列在包含 \(3/4\) 的紧集上没有共同上界，不能从点值数据抽出覆盖整个区域的收敛子列。

\ProseParagraphBreak

内部聚点的要求也不能改成仅在边界聚集，即使全族在区域上已有统一界。下面在单位圆盘内构造一个有界全纯函数，使它在无内部聚点的无限集合上为零。令
\[
 W(z)=\frac{1+z}{1-z},\quad
 q(z)=e^{-W(z)}-e^{-2}\quad(z\in\mathbb D).
\]
分母在 \(\mathbb D\) 无零，且
\[
 \operatorname{Re}W(z)=\frac{1-|z|^2}{|1-z|^2}>0,
\]
所以 \(q\) 全纯，并有 \(|q(z)|\le1+e^{-2}\)。对每个整数 \(k\ge1\)，取
\[
 z_k=\frac{1+2\pi\mathrm{i}k}{3+2\pi\mathrm{i}k}.
\]
其模平方为 \((1+4\pi^2k^2)/(9+4\pi^2k^2)<1\)，故 \(z_k\in\mathbb D\)；又有 \(z_k\rightarrow1\)，且各点互不相同。直接代入得 \(W(z_k)=2+2\pi\mathrm{i}k\)，从而 \(q(z_k)=0\)。于是 \(E=\{\,z_k\mid k\ge1\,\}\) 完全包含于定义域，却只在边界点 \(1\) 聚集。

\ProseParagraphBreak

现在令 \(f_{2n}=0\)、\(f_{2n-1}=q\)。这是一列在 \(\mathbb D\) 上共同有界的全纯函数；在每个 \(z\in E\) 上，整个数列 \(f_n(z)\) 恒为零，当然收敛。但 \(q(0)=e^{-1}-e^{-2}\ne0\)，所以 \(f_n(0)\) 在两个不同的值之间交替，整列连逐点收敛都不成立。本例中的两条子列极限 \(0\) 与 \(q\) 确实存在，也在 \(E\) 上相等；失效的是恒等定理所需的内部聚点，因而不能把子列的存在提升为整列收敛。
